\section{Causal Preservation}\label{sec:causal}

\paragraph{Where GPT-2 carries the persona effect.} Before asking whether an SAE preserves a causal effect, we
located one in the original model. For each RepE-LM pair, the honest and dishonest instructions change the next-token
distribution by a mean $\mathrm{KL}(p_H\|p_U)$ of 0.033 nats. Copying the honest activation into the dishonest run at
one layer and position group restores a share of that divergence. The effect is carried by the persona word early
(share 0.97 at the layer-1 persona slot) and by the final token late (0.95 at layer 11), with a handover at layers 7--8
(Figure~\ref{fig:m3}, left). Of 65 cells, 50 pass the pre-registered test against 200 matched-norm random vectors with
Holm correction. The analysis sites were selected on training prompts and replicated on validation prompts: the layer-1
persona slot, the layer-11 final token and, by a pre-registered amendment, the layer-8 final token. Synonym and template
controls are described in Appendix~\ref{app:causal}.

\paragraph{How much of the effect the reconstruction retains.} At a site, the activation difference between the two
prompts splits exactly into a reconstructed part and an error part. We patch each part into the dishonest run and
measure the share of the full effect it restores, together with their interaction. The reconstructed difference
restores 93--97\,\% of the effect at the layer-1 persona slot but only 62--68\,\% at the layer-8 and layer-11 final token
(Figure~\ref{fig:m3}, right). At the later sites ReLU's error part alone carries 40--45\,\%, and for TopK and JumpReLU the
interaction is large (0.26--0.38 of the effect), so their shares cannot be read additively. Higher explained variance
does not imply higher causal preservation: at layer 11 TopK has higher EV than ReLU (0.815 versus 0.757) but retains
less of the effect (0.62 versus 0.68).

\paragraph{Feature patches and their interpretation.} Patching only the top-$k$ SAE latents, with the error term kept,
restores more of the effect than 200 frequency- and magnitude-matched random latent sets for every SAE and site
($p=0.005$, the minimum attainable), and the effect is distributed: one latent carries 2--10\,\% of the reconstructed
share and 64 latents carry 47--74\,\%. Because the error term is retained, every such patch is equivalent to adding a
single vector to the original activation; these experiments therefore show which directions the SAE selects and how
much of the effect its reconstruction preserves, not that downstream computation is mediated through SAE
nonlinearities (Appendix~\ref{app:causal}).

\paragraph{A low-rank reference.} In a descriptive comparison (\evlabel{validation}), the effect is low-rank: at layer 8 a supervised rank-4 subspace carries more of it
than the whole SAE reconstruction (0.91 versus 0.62--0.67), whereas random subspaces carry at most 0.13 even at rank 64
(Appendix Table~\ref{tab:r5}). The partial preservation at later sites is thus a property of the SAE basis, not of the
task. The effect is small at layer 8 (0.022 nats), so one SAE's preservation share there can only be resolved to about
0.05--0.11, and smaller differences between SAEs at that site are not interpretable.

\begin{figure}[t]
  \centering
  \begin{minipage}[t]{0.49\linewidth}
    \centering
    \includegraphics[width=\linewidth]{fig_m3_sites}
  \end{minipage}\hfill
  \begin{minipage}[t]{0.49\linewidth}
    \centering
    \includegraphics[width=\linewidth]{fig_m4_preservation}
  \end{minipage}
  \caption{Causal site and causal preservation in GPT-2 Small (\evlabel{validation}). Left: share of the
  persona-instruction effect restored by patching one position group at one layer of the original model. Right: share
  carried by the SAE-reconstructed activation difference at the three selected sites (mean $\pm$ sd over 3 seeds;
  templates T2--T4). Per-architecture decomposition values are in Appendix Table~\ref{tab:m4}.}
  \label{fig:m3}\label{fig:m4}
\end{figure}
